\section{Related Work}
\label{sec:related_work}

\subsection{Flood Prediction}
Data-driven approximations have increasingly supplemented traditional, computationally intensive numerical hydrological models. Machine learning algorithms, particularly Random Forest, Support Vector Machines, and Artificial Neural Networks, have demonstrated strong baselines for forecasting river stages and localized flood risks \cite{Mosavi2018}. However, some conventional evaluation strategies may not reflect the temporal constraints of operational forecasting, occasionally utilizing randomized cross-validation that is susceptible to temporal leakage \cite{Kaufman2012}. This study extends these foundational approaches by strictly enforcing chronological validation and integrating flood targets into a broader multi-hazard framework.

\subsection{Heatwave and Extreme-Climate Prediction}
Predicting extreme heat typically relies on detecting anomalous spatial-temporal climatic patterns \cite{Chattopadhyay2020}. While complex deep learning models are frequently deployed for such tasks, determining accurate target boundaries remains precarious. Model inputs occasionally suffer from target-proxy leakage, where deterministic indicators (e.g., maximum daily temperatures) trivially expose the target classification to the algorithm during inference. This study explicitly mitigates this by restricting the predictive feature space, dropping direct proxy variables to ensure the model maps underlying drivers rather than tautological indicators.

\subsection{Compound Climate Events}
The characterization of concurrent and interacting environmental extremes has become a critical focus within climate risk assessment \cite{Zscheischler2018}. A rigorous typology classifies compound events into preconditioned, multivariate, spatially compounding, and temporally compounding mechanisms \cite{Zscheischler2020}. While extensive conceptual frameworks exist, operational forecasting systems frequently treat hazards independently. This framework explicitly incorporates a compound-specific predictive target, modeling the intersection of hydro-meteorological extremes to better reflect realistic concurrent risks.

\subsection{Sequence and Tree-Based Models}
Environmental time-series prediction heavily leverages both sequence and tree-based paradigms. Recurrent architectures such as Long Short-Term Memory (LSTM) and GRU \cite{Cho2014} are highly effective at capturing intrinsic temporal structures directly from sequential inputs \cite{Kratzert2018}. Conversely, Random Forest \cite{Breiman2001}, XGBoost \cite{Chen2016}, and LightGBM \cite{Ke2017} operate powerfully on tabular data by employing extensive feature engineering, such as lagged variables and rolling aggregations. In this implementation, the GRU requires a 14-day sequence and was therefore retained as an offline benchmark rather than the stateless operational predictor.

\subsection{Explainability and Calibration}
As machine learning models are inherently opaque, post-hoc interpretation is necessary for scientific trust. SHapley Additive exPlanations (SHAP) \cite{Lundberg2017} are widely utilized to characterize model-level feature dependence without strictly implying physical causality. Furthermore, because hazard events are rare, models frequently output skewed confidence scores. Techniques such as Platt scaling \cite{Platt1999} and Isotonic regression \cite{Zadrozny2002} are critical methodological steps to correct output distributions, transforming uncalibrated raw scores into reliable predictive probabilities.
